%%%%%%%%%%%%%%%%%%%%%%%%%%%%%%%%%%%%%%%%%%%%%%%%%%%%%%%%%%%%%%
%% Rozdział 5: Tłumaczenie tekstu (gloss2text / text2gloss)
%%%%%%%%%%%%%%%%%%%%%%%%%%%%%%%%%%%%%%%%%%%%%%%%%%%%%%%%%%%%%%

\section{Tłumaczenie tekstu (gloss2text / text2gloss)}
\label{chap:text-translation}

W rozdziale~\ref{chap:isolated-gloss-recognition} omówiono etap wideo--glosy
potoku przetwarzania wprowadzonego w podrozdziale~\ref{sec:motivation-context},
polegający na przekształcaniu wyodrębnionego nagrania znaku migowego w etykietę
glosy. Niniejszy rozdział dotyczy pozostałego etapu przejścia między reprezentacją
symboliczną a językiem naturalnym: tłumaczenia między sekwencjami glos a tekstem
w języku naturalnym w obu kierunkach -- glosy--tekst i tekst--glosy. Pozwala to
przekształcać rozpoznane sekwencje glos w płynne, naturalne zdania, a kierunek
odwrotny wykorzystywać do augmentacji danych i ewaluacji.

Aktualna implementacja modułu rozpoznawania w czasie rzeczywistym tworzy już
uporządkowaną sekwencję rozpoznanych etykiet glos. Sekwencja ta jest wyświetlana
w interfejsie i może zostać wyczyszczona lub skrócona o ostatnią glosę; stanowi
ona bezpośredni interfejs wejściowy dla przyszłego modułu gloss2text. W obecnym
stanie projektu nie zaimplementowano jeszcze modelu tłumaczącego tę sekwencję
na zdanie języka naturalnego. Z tego względu niniejszy rozdział opisuje
planowany etap tłumaczenia, a nie działający element potoku wnioskowania.
\todo{DO UZUPEŁNIENIA: Tematyka tego rozdziału nie została jeszcze uwzględniona
w dotychczasowych raportach z prac badawczych i programistycznych projektu.
Rozdział należy opracować po przeprowadzeniu eksperymentów z tłumaczeniem między
glosami a tekstem. Poniższe podrozdziały przedstawiają jego planowaną strukturę.}

\subsection{Podejście i dane}
\label{sec:text-translation-approach-data}

\todo{DO UZUPEŁNIENIA: Należy opisać równoległy korpus glos i tekstu (lub korpusy)
wykorzystywany do trenowania i ewaluacji, sposób pozyskiwania sekwencji glos
z wyników etapu wideo--glosy opisanego
w rozdziale~\ref{chap:isolated-gloss-recognition} lub z istniejących anotowanych
korpusów, zastosowane metody czyszczenia i normalizacji sekwencji glos oraz tekstu,
a także metodykę podziału danych na zbiory treningowy, walidacyjny i testowy.}

\subsection{Model i trenowanie}
\label{sec:text-translation-model-training}

\todo{DO UZUPEŁNIENIA: Należy opisać architekturę typu sekwencja--sekwencja
wykorzystywaną do tłumaczenia gloss2text i text2gloss (np.\ transformer
w architekturze enkoder--dekoder, dostrajanie wstępnie wytrenowanego modelu
językowego), strategię tokenizacji glos, funkcję celu uczenia, optymalizator
oraz hiperparametry. Wybory architektoniczne należy odnieść do modeli NLP
omówionych w podrozdziale~\ref{sec:nlp-models-sign-language}.}

\subsection{Ewaluacja w obu kierunkach}
\label{sec:text-translation-evaluation}

\todo{DO UZUPEŁNIENIA: Należy przedstawić wyniki ilościowe dla obu kierunków
 tłumaczenia -- glosy--tekst i tekst--glosy (np.\ BLEU, ROUGE lub inne odpowiednio
 dobrane metryki tłumaczenia maszynowego), przykłady jakościowe uzyskanych
 tłumaczeń oraz analizę typowych błędów. Ewaluacja ta zostanie ponownie
 wykorzystana do oceny etapu tłumaczenia tekstu w całościowym potoku przetwarzania,
 ocenianym w rozdziale~\ref{chap:sentence-level-translation}.}